\section*{Chapitre \ref{ch:rc:curve-construction} --- Construction de la courbe}

\begin{solution}{exo:rc:curve-construction:1}
Le forward constant égale le forward discret sur l'intervalle~: $(3.30\times3 -
3.25\times2)/(3-2) = 3.40\%$, à 2.5 ans comme partout entre les \omterm{def:rc:curve-construction:pillar}{piliers}.
\end{solution}

\begin{solution}{exo:rc:curve-construction:2}
Le taux du contrat à terme vaut $100-96.72 = 3.28~\%$~; moins l'ajustement de convexité,
$3.28-0.006 = 3.274~\%$, simple ACT/360 sur le trimestre.
\end{solution}

\begin{solution}{exo:rc:curve-construction:3}
Sa période est déjà valorisée par les dépôts et les premiers contrats à terme~: deux
instruments valoriseraient une seule période. La courbe implique un taux de swap à un an de
$3.4639\%$~; ramené à $3.4839\%$, le système reste carré (son \omterm{def:rc:curve-construction:pillar}{pilier} se situe
entre le troisième et le quatrième contrat à terme) et se calibre exactement, mais la
seule liberté restante tient aux deux semaines comprises entre le 15 et le 29 septembre 2027~:
avec des forwards constants, le forward y passe de $3.242~\%$ à $3.753~\%$, puis
retombe à $3.150~\%$ pour le reste du trimestre du quatrième contrat à terme. L'exactitude s'achète
avec un pic de 51 points de base que tout instrument portant sur cette quinzaine hérite.
\end{solution}

\begin{solution}{exo:rc:curve-construction:4}
$f = z + Tz'$. La pente gauche $(3.2883-3.4165)/(2.0027-1.2110) = -0.1619~\%$ par an~:
$f^- = 3.2883 + 2.0027\times(-0.1619) = 2.964~\%$. La pente droite
$(3.3095-3.2883)/1 = 0.0212~\%$~: $f^+ = 3.2883+2.0027\times0.0212 = 3.331~\%$. Un
saut de 37 points de base, produit par la règle.
\end{solution}

\begin{solution}{exo:rc:curve-construction:5}
$f(0,10) = \beta_0 + \beta_1e^{-10/\tau_1} + \beta_2(10/\tau_1)e^{-10/\tau_1} +
\beta_3(10/\tau_2)e^{-10/\tau_2} = 3.9407~\%$. Lorsque $T\to0$ chaque facteur de charge
$(1-e^{-x})/x\to1$ et les termes en bosse s'annulent~: $z(0)=\beta_0+\beta_1 =
2.3264~\%$.
\end{solution}

\begin{solution}{exo:rc:curve-construction:6}
Les couvertures à quatre, cinq et douze ans (et les petites à trois et quinze
ans)~: avec des forwards constants, la date à huit ans ne dépend que des \omterm{def:rc:curve-construction:pillar}{piliers} à
sept et dix ans. La spline cubique est un ajustement global~: déplacer un \omterm{def:rc:curve-construction:pillar}{pilier} fait
fléchir toute la courbe des taux zéro, si bien que le facteur d'actualisation à huit ans
réagit à des \omterm{def:rc:curve-construction:pillar}{piliers} éloignés, avec des signes alternés~; ces couvertures couvraient l'interpolation, non le marché.
\end{solution}

\begin{solution}{exo:rc:curve-construction:7}
Avec des chocs d'un point de base, les compartiments monotones convexes valent $-5\,663$ (cinq
ans), $+50\,916$ (sept), $+33\,188$ (dix) et $-7\,232$ (douze), soit au total
71\,216~USD~; avec un centième de point de base, remis à l'échelle, ils totalisent 69\,610,
et un choc parallèle d'un point de base donne 69\,579. Le cas du schéma dépend
des données~: un choc d'un point de base sur un \omterm{def:rc:curve-construction:pillar}{pilier} fait changer de cas sur ses intervalles,
et la différence finie capte plus que la pente locale~; le petit choc, lui, mesure la dérivée,
et ses compartiments s'additionnent.
\end{solution}

\begin{solution}{exo:rc:curve-construction:8}
Revaloriser les données d'entrée est nécessaire, mais non suffisant~: toute interpolation du
chapitre le fait, et elles divergent sur chaque instrument situé entre les \omterm{def:rc:curve-construction:pillar}{piliers} (de dizaines
de points de base pour les forwards, de milliers de dollars pour les swaps hors \omterm{def:rc:curve-construction:pillar}{pilier}).
La courbe doit aussi être jugée sur ses forwards (ni dents de scie, ni dépassement), sur sa
\omterm{def:rc:curve-construction:locality}{localité de l'interpolation} et sur les couvertures qu'elle implique, ainsi que sur des instruments
hors échantillon, comme les swaps hors \omterm{def:rc:curve-construction:pillar}{pilier} échangés sur le marché.
\end{solution}

\begin{solution}{pb:rc:curve-construction:1}
\textbf{1.} Chacun a été négocié au taux au pair de la courbe du vendredi.
\textbf{2.} Non~: chaque donnée se revalorise encore exactement avec la nouvelle règle, puisque
la calibration l'impose.
\textbf{3.} La nouvelle courbe donne des forwards plus élevés entre 15 et 20 ans
(et donc plus bas ailleurs dans cette plage de \omterm{def:rc:curve-construction:pillar}{piliers})~: le taux au pair à 17 ans
monte d'environ un point de base, de $3.9975~\%$ à $4.0074~\%$.
\textbf{4.} Tous les swaps hors \omterm{def:rc:curve-construction:pillar}{pilier}, soit les sept ici~; un swap arrivant à échéance
sur un \omterm{def:rc:curve-construction:pillar}{pilier} ne changerait pas de valeur.
\textbf{5.} $3.5005~\%$ puis $3.5007~\%$~: six ans se situe dans l'intervalle de cinq à sept ans,
où les deux règles donnent presque le même forward moyen.
\textbf{6.} $-\text{USD}~90\,062$.
\textbf{7.} Une perte. Le swap receveur à 17 ans perd 863\,596~USD et le swap payeur à 22 ans
gagne 678\,267~USD~; le swap payeur à 13 ans gagne 234\,745~USD~; le swap receveur à 11 ans, le swap payeur à 9 ans
et le swap receveur à 8 ans se compensent à peu près.
\textbf{8.} Il s'agit d'une variation de valorisation, non d'un résultat de trading~: le product control et le comité de valorisation
décident de la manière de présenter un changement de méthode, souvent sur une ligne distincte,
et le changement est approuvé avant sa mise en production.
\textbf{9.} «~La nouvelle interpolation, que la trésorerie utilise aussi, valorise nos swaps hors \omterm{def:rc:curve-construction:pillar}{pilier}
de 90\,000~USD de moins~; aucun marché n'a bougé.~»
\textbf{10.} Non~: les swaps à échéance de \omterm{def:rc:curve-construction:pillar}{pilier} sont valorisés par leurs propres
cotations, quelle que soit l'interpolation.
\textbf{11.} Swap payeur sur 0.3016~milliard de dollars de la position à quinze ans avant, 0.5076
milliard après~: $346\,584/1\,149\,075$ et $583\,281/1\,149\,075$ milliard.
\textbf{12.} Swap receveur sur 0.2490~milliard de dollars avant, 0.3350~milliard après.
\textbf{13.} La somme est le DV01 parallèle du portefeuille, fixé par le niveau de la courbe,
sur lequel les deux règles s'accordent~; les règles ne divergent que sur la manière dont le risque
se répartit entre les \omterm{def:rc:curve-construction:pillar}{piliers}.
\textbf{14.} Avec la méthode monotone convexe, un choc d'un point de base peut changer le cas
sur un intervalle et ajouter un morceau parasite~; un petit choc, mis à l'échelle, mesure
la dérivée (l'exercice~7).
\textbf{15.} La forme des forwards entre dix et quinze ans dépend désormais du \omterm{def:rc:curve-construction:pillar}{pilier} à douze ans
par les forwards aux \omterm{def:rc:curve-construction:pillar}{piliers} que la méthode estime~; les swaps à 11
et à 13 ans lisent cette forme.
\textbf{16.} Pour~: des forwards continus et sans dépassement, et la méthode de la courbe
officielle de la trésorerie, si bien que les valorisations s'alignent sur une référence publique.
Contre~: ses compartiments dépendent de la taille du choc, et elle est plus difficile à expliquer.
\textbf{17.} Comme changement de modèle~: documenté, validé indépendamment, son impact en P\&L
mesuré lors de calculs parallèles avant la bascule, approuvé et
présenté séparément.
\textbf{18.} Une réserve pour la dispersion des valeurs entre les interpolations acceptables
(ici au moins 90\,000~USD sur ce portefeuille), libérée au fur et à mesure
que les positions hors \omterm{def:rc:curve-construction:pillar}{pilier} du portefeuille s'éteignent.
\textbf{19.} Résultat nommé~: \emph{la bascule d'interpolation} coûte
$-\text{USD}~90\,062$ sans mouvement de marché, et la couverture à quinze ans augmente de
68~\%, passant de 0.30 à 0.51~milliard de dollars.
\textbf{20.} Tout ce qui se situe entre les cotations~: les forwards, la valeur de chaque
instrument hors \omterm{def:rc:curve-construction:pillar}{pilier}, et l'endroit où se logent ses couvertures.
\end{solution}

\begin{solution}{iq:rc:curve-construction:1}
Le bootstrap de la courbe résout un \omterm{def:rc:curve-construction:pillar}{pilier} à la fois, ce qui fonctionne lorsque chaque
instrument ne dépend que des \omterm{def:rc:curve-construction:pillar}{piliers} allant jusqu'à sa propre échéance, comme avec une
interpolation linéaire ou à forward constant. Un ajustement global résout simultanément tous les \omterm{def:rc:curve-construction:pillar}{piliers} (méthode de Newton
sur le système complet)~; il est nécessaire pour les interpolations non locales (splines,
splines sous tension, \omterm{def:rc:curve-construction:monotone}{interpolation monotone convexe} avec ses forwards aux \omterm{def:rc:curve-construction:pillar}{piliers}), pour les instruments qui
dépendent de \omterm{def:rc:curve-construction:pillar}{piliers} ultérieurs, ou pour les ajustements par moindres carrés avec plus d'instruments que
de \omterm{def:rc:curve-construction:pillar}{piliers}.
\iqlookfor{la structure triangulaire inférieure comme condition, et quand elle échoue.}
\end{solution}

\begin{solution}{iq:rc:curve-construction:2}
Le forward vaut $z+Tz'$~: la pente de la courbe des taux zéro saute à chaque \omterm{def:rc:curve-construction:pillar}{pilier}, et
le saut est multiplié par l'échéance, de sorte que la courbe à terme est en dents de scie,
avec des marches de plusieurs dizaines de points de base aux longs \omterm{def:rc:curve-construction:pillar}{piliers}, ce qui
fausse ensuite la valorisation des instruments à départ différé et des caps.
\iqlookfor{la formule $f = z + Tz'$.}
\end{solution}

\begin{solution}{iq:rc:curve-construction:3}
Ajustement exact des données d'entrée~; forwards positifs, continus et lisses~; absence de dépassement~;
localité d'un choc~; stabilité (une faible variation des données entraîne une faible variation de la courbe)~;
linéarité par rapport aux données (compartiments additifs)~; rapidité~; explication simple.
La régularité entre en conflit avec la localité (une spline est lisse et globale)~; l'absence
de dépassement entre en conflit avec la linéarité (l'\omterm{def:rc:curve-construction:monotone}{interpolation monotone convexe} est non linéaire).
\iqlookfor{une liste et au moins un véritable compromis.}
\end{solution}

\begin{solution}{iq:rc:curve-construction:4}
Une interpolation non locale, comme une spline sur les taux zéro~: le risque sur les \omterm{def:rc:curve-construction:pillar}{piliers}
éloignés provient de la règle d'interpolation de la courbe, et non des flux du swap. Vérifier
les réglages de la courbe, puis sa localité en faisant varier une donnée et en traçant la
variation du forward.
\iqlookfor{diagnostiquer l'interpolation, et non la transaction.}
\end{solution}

\begin{solution}{iq:rc:curve-construction:5}
Un passage de fin de période~: un \omterm{def:rc:curve-construction:pillar}{pilier} supplémentaire (ou un spread additif) sur les
quelques jours autour de la date, ajusté sur les contrats à terme ou dépôts qui l'encadrent,
afin que son coût ne soit pas dilué sur un trimestre. Dates de réunion~: un \omterm{def:rc:curve-construction:pillar}{pilier} par
réunion, avec des forwards constants entre eux à court terme, ajustés sur des swaps de taux
au jour le jour ou des contrats à terme, puis une règle lisse au-delà de deux ou trois ans.
\iqlookfor{structures en marches traitées explicitement, et non lissées.}
\end{solution}

\begin{solution}{iq:rc:curve-construction:6}
Pas nécessairement. Pour une courbe linéaire par rapport à ses données, le compartiment est
indépendant du choc~; pour une règle non linéaire comme l'\omterm{def:rc:curve-construction:monotone}{interpolation monotone convexe}, ou
avec une tolérance de solveur lâche, il ne l'est pas. Vérifier d'abord la tolérance du solveur
(un choc d'un dixième de point de base doit être nettement supérieur à celle-ci), puis comparer avec un petit choc agrandi~;
si la courbe est non linéaire, rapporter la dérivée obtenue avec le petit choc et l'indiquer explicitement.
\iqlookfor{la non-linéarité par rapport au bruit numérique, et la manière de les distinguer.}
\end{solution}
